\chapter{Algorithm for normalizing}\label{appendix:alg-norm}

For normalizing transformers, first we make them meet \projection\ and then \concentrated.

\section*{Algorithm for \projection}

Notice that every linear transformation $M: \Floating^d \to \Floating^{|\Sigma|}$ with $d \geq |\Sigma|$ can be decomposed into two steps. 
\footnote{We can always add more coordinates to the embedding if $d < |\Sigma|$. Since the language is fixed, it will be a constant number of coordinates, thus not changing the class of the transformer.}

\[
    \left( \begin{matrix} x_1    \\ \vdots \\ x_d    \end{matrix} \right) 
    \xrightarrow{M'}
    \left( \begin{matrix} (Mx)_1 \\ \vdots \\ (Mx)_{|\Sigma|} \\ v \end{matrix} \right)
    \xrightarrow{projection}
    \left( \begin{matrix} (Mx)_1 \\ \vdots \\ (Mx)_{|\Sigma|} \end{matrix} \right)
\]

\noindent where $v \in \Floating^{d-|\Sigma|}$ can have any values.

Let $\TF$ be a transformer and $M$ be the matrix of the $\OUTPUT$ layer.
Using the decomposition just presented, and the strategy defined in Section \ref{lin-tranf} for applying linear transformations, we can change the matrix of the $\OUTPUT$ layer to be only a projection.
The linear transformation $M$ has to be applied to the last vector, since it is the one going from the body of $\TF$ to the $\OUTPUT$ layer. 

\medskip

\section*{Algorithm for \concentrated}

Let $\TF$ be a transformer satisfying \projection, we give an algorithm for making it also meet \concentrated.

We use the maximum operation described in Section \ref{maximum} to make all the positions of the last vector $0$ except the one that was the maximum before.
Then, if the maximum is positive, multiplying by $\maxrep$, then subtracting $1$ from every coordinate and multiplying by $\maxrep$ again makes $\TF$ meet \concentrated.

\[
    \left( \begin{matrix} 0 \\ \vdots \\ 0 \\ x_{max} \\ 0 \\ \vdots \\ 0    \end{matrix} \right) 
    \xrightarrow[\maxrep]{\text{multiplying twice by }\maxrep}
    \left( \begin{matrix} 0 \\ \vdots \\ 0 \\ \maxrep \\ 0 \\ \vdots \\ 0    \end{matrix} \right) 
    \xrightarrow{\text{subtracting } 1}
    \left( \begin{matrix} -1 \\ \vdots \\ -1 \\ \maxrep -1 \\ -1 \\ \vdots \\ -1    \end{matrix} \right) 
    \xrightarrow[\maxrep]{\text{multiplying by}}
    \left( \begin{matrix} \minrep \\ \vdots \\ \minrep \\ \maxrep \\ \minrep \\ \vdots \\ \minrep    \end{matrix} \right) 
\]

It could be the case that the maximum was negative, in that case the resulting vector is not what we presented.
We show a mechanism for making the maximum strictly positive so we can insert it in the algorithm before subtracting $1$.

\medskip

At this point our vector looks like $(0, \dots, 0, x_{max}, 0, \dots, 0)$, we are going to transform it to $(0, \dots, 0, |x_{max}|, 0, \dots, 0)$.
Recall that the coordinates we are interested in are only the first $|\Sigma|$.
Let us call $v$ to the block containing these coordinates.
Since $\relu(x) + \relu(-x) = |x|$ for all $x$, our algorithm does the following:

\[
    \left(\begin{matrix}  v \\ \vdots \end{matrix}\right)
    \xrightarrow[]{}
    \left(\begin{matrix}  v \\ -v \\\vdots \end{matrix}\right)
    \xrightarrow[]{}
    \left(\begin{matrix}  \relu(v) \\ \relu(-v) \\\vdots \end{matrix}\right)
    \xrightarrow[]{}
    \left(\begin{matrix}  \relu(v) + \relu(-v) \\ \vdots \end{matrix}\right)
    = \left(\begin{matrix}  |v| \\ \vdots \end{matrix}\right)
\]

\noindent where $\relu$ is applied position wise. 


The first step is done in the same way as the first half of the linear transformations.
For the second step, notice that $\relu(x)$ can be computed with one $\FF$ layer ($W_2 = W_1 = id$, $b_2 = b_1 = 0$).
Thus, $\relu(v)$ and $\relu(-v)$ can be computed with one $\FF$ layer.
Then, in one step $\relu(v)$ and $\relu(-v)$ are computed and stored in fresh coordinates.
Next, they are moved to the desired positions.
Finally, for adding both vectors we can use one $\FF$ layer with the following parameters.
Observe that this works since all the coordinates of $\relu(-v)$ are non-negative, thus $\relu(\relu(-v)) = \relu(-v)$.

\bigrafa{Ahí mejoré el párrafo, además de que cambié la forma de computar el segundo step a algo mucho más fácil. No creo que valga la pena que expanda mucho más no?}

\begin{align*}
    & W_1 = id^{d\times d}
    & W_2  = \left(\begin{matrix}
        0^{|\Sigma| \times |\Sigma|}    & id^{|\Sigma| \times |\Sigma|}   & 0^{|\Sigma| \times (d - 2|\Sigma|)}\\
        0^{d-|\Sigma| \times |\Sigma|}  & 0^{d-|\Sigma| \times |\Sigma|}  & 0^{d-|\Sigma| \times d-2|\Sigma|} \\
    \end{matrix}\right) &
    & b_1 = b_2 = \left(\begin{matrix}
        0^{d\times d}
    \end{matrix}\right) 
\end{align*}
